% TOP_PROBLEM: {"id":30007679,"problem_number":"LOCAL-30007679","title":"Common ownership mechanisms","category_id":21,"category":{"id":21,"name":"economics","display_name":"Economics","description":"Open economic theory statements, published research agendas, and proposed scoped specifications; question provenance and historical priority are qualified per record.","slug":"economics","order_index":21,"created_at":"2026-10-08T00:00:00Z"},"status":"research_question","external_url":null,"legacy_ids":[],"published":false,"statement":"Identify price effects of institutional common ownership and its governance channels in grocery manufacturing, separating ownership shifts from common shocks.","economics":{"original_id":168,"field":"Industrial organization and competition","kind":"identification","formalization_basis":"adapted_research_question","source_status":"research_agenda","priority_status":"earliest_found","priority_note":"Earliest relevant explicit question or agenda passage located in this audit. No claim of first-ever formulation; earlier versions and later solutions were not exhaustively traced.","earliest_verified_reference":{"authors":["Carl Shapiro","Ali Yurukoglu"],"title":"Trends in Competition in the United States: What Does the Evidence Show?","year":2024,"url":"https://web.stanford.edu/~ayurukog/trendsincompetition.pdf","doi":"10.3386/w32762","locator":"Section 5, printed pp. 33–34, common-ownership discussion","evidence_summary":"Mixed price evidence motivates further work on mechanisms and outcomes."},"current_reference":{"authors":["Carl Shapiro","Ali Yurukoglu"],"title":"Trends in Competition in the United States: What Does the Evidence Show?","year":2024,"url":"https://web.stanford.edu/~ayurukog/trendsincompetition.pdf","doi":"10.3386/w32762","locator":"Section 5, printed pp. 33–34, common-ownership discussion","evidence_summary":"Mixed price evidence motivates further work on mechanisms and outcomes."},"formalization_note":"The source verifies a broader question or research agenda; the population, estimand, equations and restrictions here are our scoped adaptation. This is not a quoted canonical conjecture, and the audit does not prove contemporary unsolved status for every special case."},"statusQualification":"Published question or research agenda verified at the cited locator. The mathematical specification is adapted for this catalogue and includes additional assumptions; source publication does not establish first-ever priority or certify this exact model remains unsolved. The source verifies a broader question or research agenda; the population, estimand, equations and restrictions here are our scoped adaptation. This is not a quoted canonical conjecture, and the audit does not prove contemporary unsolved status for every special case.","statusReviewedAt":"2026-10-08"}
% FIRST_POSTED: 2026-10-08
% ENTRY_KIND: statement_only
% !TeX program = lualatex
\documentclass[11pt]{article}
\usepackage{fontspec}
\usepackage[a4paper,margin=25mm]{geometry}
\usepackage{amsmath,amssymb,mathtools}
\usepackage{xurl}
\usepackage[hidelinks]{hyperref}
\newcommand{\sref}[2]{\hyperlink{source:#1:#2}{[S#2]}}
\setlength{\emergencystretch}{3em}
\begin{document}
% problemId:30007679
\section*{Common ownership mechanisms}\label{30007679}
\subsection{Definitions and mathematical statement}
Consider publicly traded United States grocery manufacturers observed annually for ten years. Let \(O_t\) be the matrix of institutional investors' voting shares in these firms, and let \(G_t\) record investor engagements, compensation provisions and board decisions. Define \(p_j(O,G,x)\) as firm \(j\)'s equilibrium price under ownership \(O\), governance arrangement \(G\), and market fundamentals \(x\), including demand and production costs. For a prespecified ownership change from \(O^0\) to \(O^1\), target the total price effect \(E[p_j(O^1,G(O^1),x)-p_j(O^0,G(O^0),x)]\), where \(G(O)\) is governance induced by ownership and expectation averages over products and markets. Also target the controlled effect of ownership with governance held at a fixed feasible \(G^0\). Identify a source of ownership variation independent of latent competitive prospects, and specify how investor portfolio changes can alter governance rather than assuming portfolio-profit maximization by managers. Mediated effects require additional assumptions about ownership-induced confounders of governance and pricing; report ranges when these assumptions are unsupported. Compare predicted mechanisms with engagement and incentive records. The task seeks causal magnitude and transmission channels in this population, separating common ownership from direct cross-holdings and common shocks. Observed ownership concentration alone does not identify softer product-market rivalry.

\par\textbf{Formulation and scope.} The source verifies a broader question or research agenda; the population, estimand, equations and restrictions here are our scoped adaptation. This is not a quoted canonical conjecture, and the audit does not prove contemporary unsolved status for every special case.

\par\textbf{Open-question provenance.} An explicit question or research agenda was verified at the source locator. This does not by itself certify that every later version remains unresolved \sref{30007679}{1}.

\par\textbf{Historical priority.} Earliest relevant explicit question or agenda passage located in this audit. No claim of first-ever formulation; earlier versions and later solutions were not exhaustively traced.

\subsection{Short English statement}
Identify price effects of institutional common ownership and its governance channels in grocery manufacturing, separating ownership shifts from common shocks.

\subsection{Sources}
\begin{itemize}\raggedright
\item[\textnormal{[S1]}]\hypertarget{source:30007679:1}{} Carl Shapiro, Ali Yurukoglu. 2024. Trends in Competition in the United States: What Does the Evidence Show?. Section 5, printed pp. 33–34, common-ownership discussion.
\url{https://web.stanford.edu/~ayurukog/trendsincompetition.pdf}.
DOI: 10.3386/w32762.
{\small\textit{Earliest verified question/agenda reference; historical priority qualified below. Mixed price evidence motivates further work on mechanisms and outcomes.}}
\end{itemize}
\end{document}
